% ATS-friendly CV: one column, standard section names, real text (no images,
% icons, tables or text boxes), selectable and copy-pasteable in the PDF.
% Build:  cd cv && latexmk -pdf cv.tex
% CI refuses to publish while any \TODO{...} remains.
\documentclass[11pt,letterpaper]{article}
\usepackage[margin=0.7in]{geometry}
\usepackage[T1]{fontenc}
\usepackage{helvet}
\renewcommand{\familydefault}{\sfdefault}
\usepackage[utf8]{inputenc}
\input{glyphtounicode}
\pdfgentounicode=1
\usepackage{enumitem}
\usepackage[hidelinks]{hyperref}
\usepackage{titlesec}

\hypersetup{pdftitle={Ziqi Xu - CV}, pdfauthor={Ziqi Xu}}
\pagestyle{empty}
\setlength{\parindent}{0pt}
\titleformat{\section}{\large\bfseries}{}{0pt}{\MakeUppercase}[\vspace{-0.6em}\rule{\linewidth}{0.4pt}]
\titlespacing*{\section}{0pt}{0.9em}{0.5em}
\setlist[itemize]{leftmargin=1.2em, itemsep=1pt, topsep=1pt, parsep=0pt, after=\vspace{5pt}}
% no hyphenation: split words break ATS keyword matching
\hyphenpenalty=10000
\exhyphenpenalty=10000
\raggedright
\newcommand{\entry}[2]{\par\noindent\textbf{#1}\hfill #2\par}
\newcommand{\TODO}[1]{\textbf{[TODO: #1]}}

\begin{document}

{\LARGE\bfseries Ziqi Xu}\\[3pt]
Waterloo, ON, Canada \ | \ \href{mailto:z559xu@uwaterloo.ca}{z559xu@uwaterloo.ca}\\
\href{https://helioxv022.github.io}{helioxv022.github.io} \ | \ \href{https://github.com/HelioXv022}{github.com/HelioXv022} \ | \ \href{https://www.linkedin.com/in/ziqixu777}{linkedin.com/in/ziqixu777}

\section{Summary}
Physics M.Sc. (University of Waterloo, 2026) with research experience in neural quantum states and
quantum many-body numerics. I build and test research software in Python (PyTorch, NetKet, JAX, TeNPy)
for variational Monte Carlo, tensor networks and exact diagonalization. Seeking quantum software
engineering roles.

\section{Education}
\entry{M.Sc. in Physics, University of Waterloo, Waterloo, ON}{2024 -- 2026}
Associate Graduate Student, Perimeter Institute for Theoretical Physics (2024 -- 2026)\\
Supervisors: Prof. Robert Mann; co-supervisor Prof. Timothy Hsieh\\
Thesis: \textit{Conditional Information Flow and Neural Representability of One-Dimensional Rotated Cluster States}

\vspace{5pt}
\entry{B.Sc. in Theoretical Physics, Nankai University, Tianjin, China}{2020 -- 2024}
Culture and Sports Merit Scholarship (2023); University First-Class Athlete certification

\section{Research Experience}
\entry{Graduate Researcher, University of Waterloo and Perimeter Institute}{2024 -- 2026}
\begin{itemize}
  \item Studied what makes a quantum state learnable by neural networks, using 1D rotated cluster states: the
        exact MPS bond dimension is 2 at every angle, yet measurement outcomes can stay conditionally dependent
        over many sites (conditional mutual information, CMI, up to 1 bit).
  \item Designed an MPS-induced recurrent ansatz whose hidden state is the MPS bond vector, with separate even/odd
        sublattice tensors that match the state's alternating structure. With 48 real parameters it reached the
        N~=~64 ground state in all 5 random initializations (median energy-density error 1.6e-13); a
        bond-dimension-1 (product-state) control stayed at total energy error 32.
  \item Implemented complex RBM, tanh-RNN (variational Monte Carlo) and GRU models in NetKet and PyTorch and
        compared them on the same state: the RBM reached median energy-density error 6e-8 at N~=~8 but relies
        on exact summation that does not scale; the best RNN run was at total energy error 0.57 at N~=~64, still
        decreasing when a time limit stopped it; a GRU trained on 8192 bitstrings learned the Born distribution
        (fidelity 0.90) but not the phases. The results indicate that learnability depends on matching the
        model's hidden update to the state's structure rather than on hidden-state size.
  \item Wrote exact transfer-matrix code for CMI from a bond-dimension-2 MPS, validated against brute-force
        state-vector calculations (agreement to 1e-12) and the analytic AKLT result; tested in CI.
  \item Manuscript in preparation: \textit{Conditional Information Flow and Neural Representability of 1D
        Rotated Cluster States}.
\end{itemize}

\entry{Undergraduate Researcher, Computational Astrophysics}{2023 -- 2024}
Supervisor: Prof. Yu-Qing Lou, Department of Physics, Tsinghua University
\begin{itemize}
  \item Explored nonlinear solution families of the Lane--Emden equation for polytropic stars
        (n = 1, 3.5, 4, 4.5, 5) with MATLAB ODE solvers; tested mass-integral hypotheses with Simpson
        quadrature and step-size convergence checks.
  \item Later re-implemented the solver in Python (SciPy) with analytic-solution tests and GitHub Actions CI.
\end{itemize}

\entry{Undergraduate Researcher, Machine Learning for Nuclear Physics}{2023 -- 2024}
Supervisor: Prof. Jinniu Hu, Nankai University
\begin{itemize}
  \item Implemented Twin Neural Network Regression (Wetzel, \textit{Machine Learning} 115(7):163) to invert
        non-injective functions; ran numerical fits across families of non-injective equations and compared
        k-nearest-neighbor search strategies for resolving ambiguous inverses.
\end{itemize}

\entry{Project Manager, Super-Resolution Single-Molecule Imaging}{2021 -- 2023}
Prof. Leiting Pan's group, Nankai University; Tianjin Municipal Undergraduate Innovation Project
\begin{itemize}
  \item Used ImageJ and SMAP for single-molecule localization microscopy; performed cross-correlation
        analysis and drift correction across roughly 30{,}000 acquired images.
\end{itemize}

\section{Projects}
\entry{Collective AC Quantum Sensing Simulation}{2026}
\begin{itemize}
  \item Python package modelling an entangled sensor array; showed that phase disorder turns the linear growth of the
        collective signal with sensor count N into square-root growth (fitted exponent 0.495 for N~=~4 to 1024); pytest suite and CI.
\end{itemize}
\entry{Embedded Systems}{2026}
\begin{itemize}
  \item Arduino UNO R4 WiFi and ESP32 prototypes in C/C++, including an LED-matrix state machine.
  \item Assembled and programmed a Freenove 4WD smart car: GPIO/PWM motor control, sensor integration, and
        iterative hardware/software debugging.
\end{itemize}

\section{Teaching}
\entry{Teaching Assistant, PHYS 111, University of Waterloo}{Fall 2024}
Tutorials, grading and office hours.

\vspace{5pt}
\entry{Teaching Assistant, Nankai University}{2021 -- 2023}
Electromagnetism, C++ programming and computational physics.

\section{Awards and Competitions}
\begin{itemize}
  \item Mathematical Contest in Modeling (MCM/ICM), Problem A, \textit{Power Profile of a Cyclist} -- Finalist Award, 2022
  \item China Undergraduate Mathematical Contest in Modeling (CUMCM), wave-energy converter optimization -- Provincial Second Prize, 2022
  \item Shenzhen Cup Mathematical Modeling Challenge, 3D horizontal-well trajectory design -- Second Prize, 2022
  \item National College Student Statistical Modeling Competition -- Tianjin Municipal First Prize, 2022
  \item CUPT, team member responsible for \textit{Saving Honey} and \textit{Water Spiral} -- National Second Prize, 2022
\end{itemize}

\section{Skills}
\textbf{Programming:} Python, C++, MATLAB, LaTeX; Git, GitHub Actions, Linux, Jupyter\\
\textbf{Scientific computing:} NumPy, SciPy, PyTorch, JAX, NetKet, TeNPy, PennyLane, Matplotlib, pytest\\
\textbf{Methods:} variational Monte Carlo, neural quantum states (RBM, RNN), tensor networks (MPS, DMRG),
exact diagonalization, ODE integration, numerical quadrature\\
\textbf{Hardware:} Arduino, ESP32\\
\textbf{Languages:} English, Mandarin (native)

\end{document}
